\documentclass[11pt]{article}

\usepackage[margin=1in]{geometry}
\usepackage{amsmath,amssymb,amsthm,booktabs,hyperref}
\usepackage[T1]{fontenc}
\hypersetup{hidelinks}

\newtheorem{theorem}{Theorem}
\newtheorem{lemma}{Lemma}
\newtheorem{proposition}{Proposition}
\newtheorem{remark}{Remark}

\title{Two-scale far-pair constraints for hypothetical\\
41-point kissing codes in dimension five}
\author{Unsolved Labs}
\date{Research release R008}

\begin{document}
\maketitle

\begin{abstract}
We prove exact structural constraints on any hypothetical spherical code
$C\subset S^4$ of cardinality $41$ with pairwise inner products at most
$1/2$. Two rational Gegenbauer certificates force at least two unordered
pairs with inner product below $-442/625$ and at least ten below $-133/200$.
A third certificate bounds the degree of the moderate far-pair graph by nine.
Together with elementary graph and Euclidean geometry, these facts yield a
vertex-disjoint strong/moderate pair, at least seven moderate-edge endpoints
spanning dimension at least three, and a strict $779/1000$ bound on the
absolute inner product of suitable normalized difference axes. The
polynomial sign claims are replayed exactly by two independent Sturm
implementations, one using SymPy and one using only Python's standard-library
rational arithmetic. The result is structural and does not prove that the
five-dimensional kissing number is $40$.
\end{abstract}

\paragraph{Research provenance.}
This manuscript documents an Unsolved Labs research release generated with
frontier AI. The final claim is supported by the public proof and exact
certificate replay described below. Independent specialist review remains
pending; no conventional human authorship or external endorsement is implied.

\section{Introduction}

The kissing number $\tau_d$ is the largest number of unit spheres in
$\mathbb R^d$ that can touch a central unit sphere without overlapping.
Equivalently, it is the maximum cardinality of a spherical code
$C\subset S^{d-1}$ whose distinct inner products are at most $1/2$.
The five-dimensional problem is unresolved. Cohn and Rajagopal
\cite{cohnrajagopal2026} state that $\tau_5$ appears to be $40$, while the
best proved upper bound is $44$, and they exhibit a fourth nonisometric
40-point configuration. The upper bound $44$ comes from semidefinite
programming computations of Mittelmann and Vallentin
\cite{mittelmannvallentin2010}.

Thus any hypothetical counterexample to $\tau_5=40$ must already contain at
least $41$ points. This paper does not exclude such a code. Instead, it
shows that any 41-point code is forced to contain a substantial pattern of
very negative inner products at two distinct scales. The argument uses the
classical positive-definite Gegenbauer/Delsarte framework for spherical codes,
together with exact rational polynomial certificates. The semidefinite
programming framework for spherical codes is established prior work; see,
for example, Bachoc and Vallentin \cite{bachocvallentin2008}.

The contribution is the particular two-scale certificate and the structural
consequences extracted from it. The certificate search is not trusted in the
final proof: only the frozen rational coefficients and exact sign replay are
used.

\section{Statement}

For $a>0$, let $E_a(C)$ denote the number of unordered pairs
$\{x,y\}\subset C$ with $x\cdot y<-a$.

\begin{theorem}\label{thm:main}
Let $C\subset S^4$ satisfy $|C|=41$ and
\[
x\cdot y\le \frac12 \qquad (x\ne y).
\]
Then:
\begin{enumerate}
\item $E_{442/625}(C)\ge2$;
\item $E_{133/200}(C)\ge10$;
\item the graph $G_m$ whose edges are the pairs with inner product below
      $-133/200$ is triangle-free and has maximum degree at most $9$;
\item at least seven vertices are incident with edges of $G_m$, and those
      vertices span a subspace of dimension at least $3$;
\item there exist four distinct points $x_+,x_-,y_+,y_-\in C$ such that
\[
x_+\cdot x_-<-\frac{442}{625},\qquad
y_+\cdot y_-<-\frac{133}{200},
\]
and if
\[
e=\frac{x_+-x_-}{\|x_+-x_-\|},\qquad
f=\frac{y_+-y_-}{\|y_+-y_-\|},
\]
then
\[
|e\cdot f|<\frac{779}{1000}.
\]
\end{enumerate}
\end{theorem}

\begin{remark}
The theorem does not prove $\tau_5=40$, does not improve the proved upper
bound $44$, and makes no optimality claim for the thresholds $442/625$,
$133/200$, or $779/1000$.
\end{remark}

\section{Proof overview}

The proof has four layers.

First, an edge-count variant of the Delsarte linear-programming argument
turns a positive Gegenbauer expansion with two sign conditions into a lower
bound for the number of pairs below a chosen negative threshold.

Second, two explicit rational degree-nine certificates force at least two
``strong'' far pairs below $-442/625$ and at least ten ``moderate'' far pairs
below $-133/200$.

Third, projecting all moderate neighbors of one code point into its
orthogonal $S^3$ reduces the local degree problem to an ordinary spherical
code bound. A third exact certificate shows that at most nine moderate
neighbors are possible.

Finally, elementary graph theory and Euclidean geometry convert these
numerical facts into a disjoint-pair statement, endpoint/span constraints,
and the difference-axis separation bound.

Every polynomial sign assertion is certified by exact Sturm root counting
over $\mathbb Q$ in two independent implementations.

\section{An edge-count Delsarte lemma}

Let $P_k^{(5)}$ be the normalized Gegenbauer polynomials for $S^4$, with
$P_k^{(5)}(1)=1$.

\begin{lemma}\label{lem:edgecount}
Suppose
\[
f(t)=\sum_{k\ge0}a_kP_k^{(5)}(t),\qquad a_k\ge0,
\]
and, for some $a\in(0,1)$,
\[
f(t)\le0\quad(-a\le t\le1/2),\qquad
f(t)\le1\quad(-1\le t<-a).
\]
Then every $N$-point spherical code in $S^4$ with pairwise inner products at
most $1/2$ satisfies
\[
E_a(C)\ge
\left\lceil\frac{a_0N^2-Nf(1)}2\right\rceil.
\]
\end{lemma}

\begin{proof}
Gegenbauer positivity gives
\[
a_0N^2\le \sum_{x,y\in C}f(x\cdot y).
\]
The diagonal contributes $Nf(1)$. Each unordered pair below $-a$ contributes
at most $2$ to the ordered-pair sum, while every other off-diagonal pair
contributes at most $0$. Hence
\[
a_0N^2\le Nf(1)+2E_a(C),
\]
which proves the claim.
\end{proof}

\section{Two exact far-pair certificates}

\subsection{Strong threshold}

Define
\[
f_s(t)=\sum_k a_kP_k^{(5)}(t)
\]
using the nonzero coefficients in Table~\ref{tab:strong}.

\begin{table}[ht]
\centering
\begin{tabular}{crrrrrr}
\toprule
$k$ & 0 & 1 & 2 & 3 & 4 & 9\\
\midrule
$a_k$
& $\frac{13939907}{10^8}$
& $\frac{12103401}{25000000}$
& $\frac{139479103}{10^8}$
& $\frac{30090377}{25000000}$
& $\frac{41787773}{25000000}$
& $\frac{38538027}{50000000}$\\
\bottomrule
\end{tabular}
\caption{Strong-edge certificate.}
\label{tab:strong}
\end{table}

Exact Sturm replay proves
\[
f_s(t)<0\quad\text{for }-\frac{442}{625}\le t\le\frac12,
\]
and
\[
f_s(t)<1\quad\text{for }-1\le t\le-\frac{442}{625}.
\]
For $N=41$,
\[
N^2a_0-Nf_s(1)=\frac{209711679}{10^8}=2.09711679.
\]
Lemma~\ref{lem:edgecount} therefore gives
\begin{equation}\label{eq:strong}
E_{442/625}(C)\ge2.
\end{equation}

\subsection{Moderate threshold}

Define
\[
f_m(t)=\sum_k a_kP_k^{(5)}(t)
\]
using Table~\ref{tab:moderate}.

\begin{table}[ht]
\centering
\begin{tabular}{crrrrrr}
\toprule
$k$ & 0 & 1 & 2 & 3 & 4 & 9\\
\midrule
$a_k$
& $\frac{6113963}{50000000}$
& $\frac{3549967}{10^7}$
& $\frac{59464043}{50000000}$
& $\frac{11303433}{12500000}$
& $\frac{17480877}{12500000}$
& $\frac{29954461}{50000000}$\\
\bottomrule
\end{tabular}
\caption{Moderate-edge certificate.}
\label{tab:moderate}
\end{table}

Exact replay proves
\[
f_m(t)<0\quad\text{for }-\frac{133}{200}\le t\le\frac12,
\]
and
\[
f_m(t)<1\quad\text{for }-1\le t\le-\frac{133}{200}.
\]
Moreover,
\[
N^2a_0-Nf_m(1)=\frac{912370581}{50000000}=18.24741162.
\]
Thus
\begin{equation}\label{eq:moderate}
E_{133/200}(C)\ge10.
\end{equation}

\section{The moderate far-pair graph}

Let $G_m$ have vertex set $C$ and join $x$ to $y$ exactly when
$x\cdot y<-133/200$.

\subsection{Maximum degree}

Fix $x\in C$. For each neighbor $y_i$, write
\[
y_i=-p_ix+\sqrt{1-p_i^2}\,u_i,
\qquad
p_i>\frac{133}{200},
\qquad
u_i\in S^3\subset x^\perp.
\]
For distinct neighbors,
\[
u_i\cdot u_j\le
\frac{\frac12-p_ip_j}
{\sqrt{(1-p_i^2)(1-p_j^2)}}.
\]
For $p,q>1/2$, the right-hand side is decreasing in either variable: its
partial derivative with respect to $p$ has the sign of $p/2-q$.
Consequently
\[
u_i\cdot u_j<
c:=
\frac{\frac12-(133/200)^2}{1-(133/200)^2}
=\frac{2311}{22311}.
\]

Now use normalized Gegenbauer polynomials $P_k^{(4)}$ for $S^3$ and the
coefficients in Table~\ref{tab:local}.

\begin{table}[ht]
\centering
\begin{tabular}{crrrrr}
\toprule
$k$ & 0 & 1 & 2 & 3 & 6\\
\midrule
$b_k$
& $\frac{99999}{100000}$
& $\frac{153764653}{50000000}$
& $\frac{89219103}{25000000}$
& $\frac{163063619}{10^8}$
& $\frac{13716513}{10^8}$\\
\bottomrule
\end{tabular}
\caption{Local projected $S^3$ certificate.}
\label{tab:local}
\end{table}

The resulting polynomial $g$ is exactly negative on
$[-1,2311/22311]$, while
\[
\frac{g(1)}{b_0}
=\frac{18823697}{1999980}
=9.4119426\ldots<10.
\]
The ordinary Delsarte bound therefore implies
\begin{equation}\label{eq:degree}
\Delta(G_m)\le9.
\end{equation}

\subsection{Triangle-freeness and a matching}

If three code vectors were pairwise adjacent in $G_m$, then
\[
\|x+y+z\|^2
<
3-6\left(\frac{133}{200}\right)<0,
\]
which is impossible. Thus $G_m$ is triangle-free.

A triangle-free graph in which every two edges intersect is a star: if $xy$
and $xz$ are edges, an edge avoiding $x$ would have to be $yz$, creating a
triangle. Therefore a triangle-free graph with no matching of size two is a
star. But \eqref{eq:moderate} gives at least ten edges while
\eqref{eq:degree} bounds the degree of any star by nine. Hence $G_m$
contains two vertex-disjoint edges.

Let $G_s\subset G_m$ be the strong graph at threshold $-442/625$.
Equation~\eqref{eq:strong} gives at least two strong edges. If two strong
edges are disjoint, we are done. Otherwise take intersecting strong edges
$xy,xz$ and a two-edge matching of $G_m$. If one matching edge avoids $xy$,
it is disjoint from a strong edge. Otherwise the two matching edges must meet
$x$ and $y$ separately; the one through $y$ cannot contain $z$ because that
would create the triangle $xyz$. It is therefore disjoint from $xz$.
Thus there is always a strong edge disjoint from a moderate edge.

\section{Difference-axis separation}

Write the selected pairs as
\[
x_\pm=m\pm p,\qquad y_\pm=n\pm q.
\]
Since the endpoints are unit vectors, $m\perp p$ and $n\perp q$.
The pair thresholds imply
\[
\|m\|^2<\frac{183}{1250},\quad
\|p\|^2>\frac{1067}{1250},\quad
\|n\|^2<\frac{67}{400},\quad
\|q\|^2>\frac{333}{400}.
\]
All four cross pairs are distinct code points, so their inner products are at
most $1/2$. Adding opposite cross inequalities yields
\[
m\cdot n+p\cdot q\le\frac12,\qquad
m\cdot n-p\cdot q\le\frac12.
\]
Hence
\[
|p\cdot q|
\le\frac12-m\cdot n
\le\frac12+\|m\|\,\|n\|.
\]
For $e=p/\|p\|$ and $f=q/\|q\|$,
\[
|e\cdot f|
<
\frac{1+\sqrt{(1-442/625)(1-133/200)}}
{\sqrt{(1+442/625)(1+133/200)}}.
\]
The exact certificate verifies, without relying on the decimal display, that
the right-hand side is strictly below $779/1000$.

\section{Endpoint count and span}

Let $W$ be the vertices incident with $G_m$. The induced graph on $W$ is
triangle-free and has at least ten edges. Mantel's theorem gives at most
$\lfloor 6^2/4\rfloor=9$ edges on six vertices, so $|W|\ge7$.

If $W$ were contained in a subspace of dimension at most two, its unit
vectors would lie on a circle (or a smaller sphere) with pairwise angular
separation at least $60^\circ$. Such a circle code has at most six points.
Therefore
\[
\dim\operatorname{span}(W)\ge3.
\]
This completes the proof of Theorem~\ref{thm:main}.

\section{Exact verification and trust boundary}

The repository contains two independent exact replay paths.

The primary checker uses SymPy exact rational arithmetic and SymPy's Sturm
root counting. The independent checker uses only Python's
\texttt{fractions.Fraction}, constructs Gegenbauer polynomials from the
three-term recurrence, implements polynomial Euclidean division, and applies
Sturm's theorem directly. Both independently reproduce all sign intervals,
the two Delsarte deficits, the local $S^3$ objective, and the exact radical
comparison.

The search process that discovered the certificates is outside the final
trust boundary. The final proof depends only on the printed rational
coefficients, standard Gegenbauer positivity/Delsarte theory, elementary
graph and Euclidean deductions, and the exact replay.

The graph/geometric deductions are not currently Lean-formalized. The
repository therefore describes the result as exact computer-assisted
verification, not Lean verification. A useful future formalization would
target Lemma~\ref{lem:edgecount}, the projection reduction, the graph lemmas,
and the interface from exact Sturm certificates to polynomial sign claims.

\section{Related work and scope}

The positive-definite and semidefinite-programming framework for spherical
codes is prior work \cite{bachocvallentin2008}. Cohn and Rajagopal
\cite{cohnrajagopal2026} provide the current five-dimensional construction
context and cite the proved upper bound $44$ from
\cite{mittelmannvallentin2010}.

A targeted public-source search refreshed on 18 August 2026 did not locate a
prior source stating the exact thresholds $-442/625$ and $-133/200$ together
with the corresponding 2/10 pair-count bounds and the released
disjoint-pair, endpoint, span, and axis consequences for every hypothetical
41-point code. This is deliberately phrased as a scoped source audit rather
than an exhaustive priority claim.

\begin{thebibliography}{9}

\bibitem{cohnrajagopal2026}
Henry Cohn and Isaac Rajagopal.
\newblock Variations on five-dimensional sphere packings.
\newblock \emph{Discrete \& Computational Geometry}, 2026.
\newblock arXiv:2412.00937v3; doi:10.1007/s00454-026-00841-x.

\bibitem{mittelmannvallentin2010}
Hans D. Mittelmann and Frank Vallentin.
\newblock High accuracy semidefinite programming bounds for kissing numbers.
\newblock \emph{Experimental Mathematics}, 19(2):175--179, 2010.
\newblock arXiv:0902.1105.

\bibitem{bachocvallentin2008}
Christine Bachoc and Frank Vallentin.
\newblock New upper bounds for kissing numbers from semidefinite programming.
\newblock \emph{Journal of the American Mathematical Society},
21(3):909--924, 2008.
\newblock arXiv:math/0608426; doi:10.1090/S0894-0347-07-00589-9.

\end{thebibliography}

\end{document}